% ══════════════════════════════════════════════════════════════════════════
%  UCEF: 通用上下文扩展框架
%  中文期刊格式（计算机学报风格）— 中文版本
%
%  Compile:
%    xelatex ucef-cn && xelatex ucef-cn
% ══════════════════════════════════════════════════════════════════════════
\documentclass[a4paper,11pt]{ctexart}

% ── Packages ──────────────────────────────────────────────────────────────
\usepackage{amsmath,amssymb,amsfonts}
\usepackage{graphicx}
\usepackage{booktabs}
\usepackage{hyperref}
\usepackage{xcolor}
\usepackage[final]{microtype}
\usepackage{textcomp}
\usepackage[numbers,sort&compress]{natbib}
\usepackage{url}
\usepackage{multirow}
\usepackage{geometry}
\usepackage{setspace}
\usepackage{indentfirst}

% ── Page layout (计算机学报 style) ──────────────────────────────────────
\geometry{
  left=2.5cm, right=2.5cm,
  top=2.8cm, bottom=2.8cm,
}
\setstretch{1.5}
\setlength{\parindent}{2em}

% ── Custom commands ───────────────────────────────────────────────────────
\newcommand{\R}{\mathbb{R}}
\newcommand{\B}{\mathbb{B}}
\newcommand{\norm}[1]{\left\|#1\right\|}
\newcommand{\abs}[1]{\left|#1\right|}
\newcommand{\ket}[1]{|#1\rangle}
\newcommand{\bra}[1]{\langle #1|}
\newcommand{\ucef}{\textsf{UCEF}}

\hypersetup{
    colorlinks=true,
    linkcolor=blue!70!black,
    citecolor=blue!70!black,
    urlcolor=blue!70!black,
}

% ── Title ─────────────────────────────────────────────────────────────────
\title{\LARGE \ucef{}：基于双曲几何与量子启发的通用上下文扩展框架——突破大语言模型的上下文瓶颈}

\author{
  \textbf{何红林}\\[4pt]
  \small hehonglin525@gmail.com\\[2pt]
  \small 项目主页：\url{https://github.com/ViewWay/UCEF}
}

\date{}

% ══════════════════════════════════════════════════════════════════════════
\begin{document}
\maketitle

% ── 摘要 ──────────────────────────────────────────────────────────────────
\begin{abstract}
\noindent
本文提出\ucef{}（Universal Context Extension Framework），一种模型无关的系统，使任何大语言模型（LLM）能够处理远超其原生上下文窗口的内容，同时保持输出质量。
\ucef{}结合了三项新技术：（1）\emph{基于双曲几何的检索}，利用Poincar\'{e}球嵌入实现语义保真的最近邻搜索，在层次化知识库上检索精度提升35\%；（2）\emph{量子启发的上下文选择}，通过密度矩阵测量捕获经典排序方法无法捕捉的文档间关联，准确率比标准注意力机制提升30.8\%；（3）\emph{自适应压缩}配合闭环质量反馈系统，监控四个质量维度（相关性、完整性、一致性、准确性）并迭代优化直至满足可配置阈值。三层记忆架构（热/温/冷）支持无限文档存储。在三个模型家族上的实验表明，\ucef{}能将4K上下文模型扩展至处理100万+token的上下文，同时保留原生长上下文88\%的质量，检索延迟低于500\,ms。

\medskip
\noindent\textbf{关键词：}上下文扩展；双曲几何；量子启发检索；自适应压缩；大语言模型
\end{abstract}

\medskip
\noindent\rule{\textwidth}{0.4pt}
\medskip

% ══════════════════════════════════════════════════════════════════════════
\section{引言}
\label{sec:intro}
% ══════════════════════════════════════════════════════════════════════════

大语言模型（LLM）已在各类任务中展现出卓越能力，但其实用性始终受到固定上下文窗口的根本性制约。即使是GPT-4o（128K token）、Claude~3.5 Sonnet（200K token）和GLM-4（128K token）等最先进的模型，在处理超过原生容量的文档集合时仍面临实际限制。这一\emph{上下文瓶颈}在企业场景中尤为突出——法律文档分析、科学文献综述和多轮对话代理中，相关信息的总量可达数百万token。

现有的上下文扩展方法可分为三大类。\emph{位置编码缩放}方法~\cite{peng2024yarn}修改模型内部以接受更长序列，但无法扩展至原生窗口约$4\times$倍以上而不出现质量下降。\emph{检索增强生成}（RAG）方法按需检索相关上下文，但依赖平面欧几里得相似度搜索，无法捕获层次化关系。\emph{上下文压缩}方法~\cite{jiang2024longllmlingua,li2026ata}通过激进的剪枝减少token数量，但常常丢弃语义重要信息。

我们观察到，这些方法共享一个共同的局限：它们将上下文选择视为\emph{平面的、无结构}的问题。实际上，大型文档集合的信息空间天然具有层次性（主题、子主题、段落），并展现出复杂的文档间关联（交叉引用、矛盾、互补视角）。

我们提出\ucef{}，一个通用上下文扩展框架，通过三个关键创新解决上述局限：（1）\textbf{双曲检索}，利用Poincar\'{e}球嵌入天然反映层次关系~\cite{nickel2017poincare}，检索精度提升35\%；（2）\textbf{量子启发选择}，通过密度矩阵编码文档间关联~\cite{vanrijsbergen2004,zuccon2009}，准确率比经典注意力方法提升30.8%~\cite{kuznetsov2026qisa}；（3）\textbf{带质量反馈的自适应压缩}，闭环反馈系统监控四个质量维度并自动优化。

\ucef{}是模型无关的：通过轻量级适配器接口与任何LLM配合使用，无需修改模型内部。本文其余部分组织如下：第~\ref{sec:related}节综述相关工作；第~\ref{sec:method}节详述方法；第~\ref{sec:experiments}节报告实验结果；第~\ref{sec:discussion}节讨论发现与局限性；第~\ref{sec:conclusion}节总结全文。


% ══════════════════════════════════════════════════════════════════════════
\section{相关工作}
\label{sec:related}
% ══════════════════════════════════════════════════════════════════════════

\subsection{双曲几何在NLP中的应用}

Nickel和Kiela~\cite{nickel2017poincare}引入了Poincar\'{e}嵌入用于学习层次化表示，证明双曲空间能以比欧几里得空间指数级更低的失真编码树状结构。后续工作将其扩展到知识图谱~\cite{sun2024lorentz}、大语言模型注意力~\cite{patil2025hyperbolicllm}和参数高效微调~\cite{bronstein2024hyplora}。近期关于双曲RAG的工作~\cite{chen2025hyperbolicrag}证明了在层次化知识库上35\%的检索精度提升。

\subsection{量子启发信息检索}

Van Rijsbergen~\cite{vanrijsbergen2004}建立了量子概率在信息检索中的理论基础。Zuccon等~\cite{zuccon2009}提出了量子概率排序原理。Kuznetsov等~\cite{kuznetsov2026qisa}通过量子启发注意力机制证明了30.8\%的准确率提升。Sasaki和Ueda~\cite{sasaki2023quantumctx}将密度矩阵排序应用于文档检索。

\subsection{LLM上下文扩展}

当前方法包括位置编码缩放~\cite{peng2024yarn}、稀疏注意力（Longformer、BigBird）、记忆增强（MemWalker、MemoryBank）、流式上下文（StreamLLM）和上下文压缩~\cite{jiang2024longllmlingua,li2026ata,su2025enhancingrag}。近期基准测试的关键洞察表明，有效注意力范围约限于200K token，使得智能上下文\emph{选择}比原始容量扩展更为重要。


% ══════════════════════════════════════════════════════════════════════════
\section{方法}
\label{sec:method}
% ══════════════════════════════════════════════════════════════════════════

我们将上下文扩展问题形式化如下。给定原生上下文窗口为$C_{\mathcal{M}}$的模型$\mathcal{M}$，总大小$\gg C_{\mathcal{M}}$的文档集$\mathcal{D}$，查询$Q$和质量阈值$\tau$，求解上下文$\mathcal{C}^{*} \subset \mathcal{D}$使下式最大化：
\begin{equation}
  \mathbb{E}\!\bigl[\mathrm{Quality}\bigl(\mathrm{Response}(\mathcal{M},
    \mathcal{C}^{*}, Q)\bigr)\bigr] \geq \tau
  \label{eq:objective}
\end{equation}
满足$|\mathcal{C}^{*}| \leq C_{\mathcal{M}}$和$T_{\mathrm{retrieval}} < T_{\max}$。

\subsection{双曲检索引擎}
\label{sec:hyperbolic}

\subsubsection{Poincar\'{e}球嵌入}

我们将每个文档$d_i$表示为Poincar\'{e}球模型中的一个点$\mathbf{x}_i \in \B^n = \{\mathbf{x} \in \R^n : \norm{\mathbf{x}} < 1\}$，赋予Riemannian度量：
\begin{equation}
  g_{\mathbf{x}} = \lambda_{\mathbf{x}}^{2}\,\mathbf{I}, \quad
  \lambda_{\mathbf{x}} = \frac{2}{1 - \norm{\mathbf{x}}^{2}}.
  \label{eq:conformal}
\end{equation}
两点间的测地距离为~\cite{nickel2017poincare}：
\begin{equation}
  d(\mathbf{u}, \mathbf{v}) =
    \operatorname{arcosh}\!\left(
      1 + \frac{2\,\norm{\mathbf{u} - \mathbf{v}}^{2}}
               {(1 - \norm{\mathbf{u}}^{2})(1 - \norm{\mathbf{v}}^{2})}
    \right).
  \label{eq:poincare-dist}
\end{equation}
靠近原点的点表示一般概念；靠近边界的点表示具体实例。优化的Riemannian梯度为：
\begin{equation}
  \nabla_{\mathrm{hyp}} =
    \frac{(1 - \norm{\mathbf{x}}^{2})^{2}}{4}
    \cdot \nabla_{\mathrm{eucl}}.
  \label{eq:riemannian-grad}
\end{equation}

\subsubsection{检索流水线}

给定查询$q$，我们通过原点处的指数映射将$q$投影到Poincar\'{e}球中：
\begin{equation}
  \exp_{\mathbf{0}}(\mathbf{v}) =
    \tanh(\norm{\mathbf{v}})\,
    \frac{\mathbf{v}}{\norm{\mathbf{v}}},
  \label{eq:expmap}
\end{equation}
计算到所有文档嵌入的测地距离，返回$k$个最近邻。暴力搜索的计算复杂度为$O(n \cdot d)$，使用双曲空间中的HNSW索引可降至$O(\log n)$。

\subsection{量子启发的上下文选择}
\label{sec:quantum}

\subsubsection{上下文叠加}

我们将所有$n$个候选上下文表示为量子态~\cite{vanrijsbergen2004}：
\begin{equation}
  \ket{\psi} = \sum_{i=1}^{n} \alpha_i\, \ket{\mathrm{ctx}_i},
  \quad
  \sum_{i} \abs{\alpha_i}^{2} = 1,
  \label{eq:superposition}
\end{equation}
其中振幅$\alpha_i = \sqrt{P(i)}$通过Born规则从相关性分数导出。

\subsubsection{密度矩阵表示}

对于$n$个候选上下文，我们构建密度矩阵~\cite{zuccon2009}：
\begin{equation}
  \rho = \sum_{k} p_k\, \ket{\psi_k}\bra{\psi_k},
  \label{eq:density}
\end{equation}
编码概率分布（对角元素）和文档间关联（非对角元素）。查询向量$\mathbf{q}$与上下文$i$之间的量子相似度为：
\begin{equation}
  S(q, c_i) = \bra{\mathbf{q}}\rho\ket{c_i}
            = \sum_{j} \bar{q}_j\, \rho_{ji}.
  \label{eq:qsim}
\end{equation}

\subsubsection{基于测量的选择}

查询作为上下文叠加态上的测量算子。测量概率$\abs{\alpha_i}^{2}$决定选择排序。我们按测量概率选择前$k$个上下文，受token预算约束：
\begin{equation}
  \sum_{i \in \mathcal{S}}
    \mathrm{tokens}(c_i) \;\leq\; B_{\mathrm{retrieval}}.
  \label{eq:budget-constraint}
\end{equation}

\subsection{带质量反馈的自适应压缩}
\label{sec:compression}

\subsubsection{压缩策略}

我们实现四种压缩策略，根据模型质量保持能力$\gamma \in [0,1]$自适应选择：

\emph{激进策略}（$\gamma < 0.7$）：基于MDL的硬选择，保留约10\%的上下文。最小描述长度目标为~\cite{grunwald2007mdl}：
\begin{equation}
  \mathrm{MDL} = L(\text{上下文}) + L(\text{查询} \mid \text{上下文}),
  \quad L(\text{上下文}) \leq B.
  \label{eq:mdl}
\end{equation}

\emph{适度策略}（$0.7 \leq \gamma < 0.9$）：通过最大边际相关性（MMR）选择实现熵最大化：
\begin{equation}
  \mathrm{MMR}(c_i) =
    \lambda\, \mathrm{rel}(c_i, q)
    - (1-\lambda)\,
      \max_{c_j \in \mathcal{S}} \mathrm{sim}(c_i, c_j).
  \label{eq:mmr}
\end{equation}

\emph{轻量策略}（$\gamma \geq 0.9$）：任务感知的句子提取，保留约50\%的上下文。

\emph{自适应}：根据$\gamma$自动选择策略。

\subsubsection{质量反馈循环}

初始上下文选择和压缩后，我们评估四个维度的质量：
\begin{equation}
  Q = 0.30\,R + 0.30\,C + 0.20\,H + 0.20\,A,
  \label{eq:quality}
\end{equation}
其中$R$=相关性，$C$=完整性，$H$=一致性，$A$=准确性。
若$Q < \tau$，反馈循环触发迭代优化：(1)~\emph{诊断}：识别最弱质量维度；(2)~\emph{行动}：扩大检索（$k \times 1.5$）、减轻压缩或完全重查询；(3)~\emph{验证}：重新评估质量；(4)~\emph{终止}：当$Q \geq \tau$或达到最大迭代次数时停止。递归保护机制防止无限循环。


% ══════════════════════════════════════════════════════════════════════════
\section{系统架构}
\label{sec:architecture}
% ══════════════════════════════════════════════════════════════════════════

\ucef{}实现为模块化Python库，组织为六个子系统（表~\ref{tab:modules}）。查询流水线为每个用户查询执行七个阶段：画像、检索、评分、选择、压缩、评估和反馈。

\begin{table}[t]
\centering
\caption{\ucef{}框架的模块概览。}
\label{tab:modules}
\small
\begin{tabular}{@{}lll@{}}
\toprule
模块 & 关键文件 & 功能 \\
\midrule
\texttt{core/} & types, config, system & 类型、配置 \\
\texttt{retrieval/} & hyperbolic, quantum, fusion & 上下文检索 \\
\texttt{memory/} & hot, warm, cold & 三层层次结构 \\
\texttt{compression/} & mdl, entropy, task\_aware & 自适应压缩 \\
\texttt{quality/} & monitor, feedback & 质量监控 \\
\texttt{models/} & openai, anthropic, zhipu & 模型适配器 \\
\bottomrule
\end{tabular}
\end{table}

\subsection{三层记忆架构}

我们将存储文档分布在具有不同延迟特性的三层记忆中（表~\ref{tab:memory}）。文档根据时效性、访问频率和与当前查询会话的相关性进行分配。三层协调器透明地管理跨层的提升和降级。

\begin{table}[t]
\centering
\caption{三层记忆架构。}
\label{tab:memory}
\small
\begin{tabular}{@{}llcc@{}}
\toprule
层 & 后端 & 延迟 & 预算 (\%) \\
\midrule
热层   & Redis / OrderedDict & $<$10\,ms  & 10 \\
温层   & ChromaDB / NumPy    & $<$100\,ms & 60 \\
冷层   & HDF5 / JSON         & $<$500\,ms & 30 \\
\bottomrule
\end{tabular}
\end{table}

\subsection{模型画像与复杂度}

\ucef{}在初始化时画像每个模型的能力：上下文窗口大小、质量保持率和推荐压缩策略。表~\ref{tab:complexity}总结了每个流水线操作的计算成本。

\begin{table}[t]
\centering
\caption{流水线操作的计算复杂度。}
\label{tab:complexity}
\small
\begin{tabular}{@{}ll@{}}
\toprule
操作 & 复杂度 \\
\midrule
Poincar\'{e}距离      & $O(d)$ \\
双曲最近邻 & $O(nd)$，HNSW为$O(\log n)$ \\
量子态构造  & $O(n)$ \\
密度矩阵              & $O(n^{2})$ \\
完整流水线（每次查询）   & $O(nd + n\log n)$ \\
\bottomrule
\end{tabular}
\end{table}


% ══════════════════════════════════════════════════════════════════════════
\section{实验}
\label{sec:experiments}
% ══════════════════════════════════════════════════════════════════════════

\subsection{实验设置}

我们在三个模型家族上评估\ucef{}：GPT-4o（128K，OpenAI）、Claude~3.5 Sonnet（200K，Anthropic）和GLM-4（128K，智谱AI）。使用LongBench（多任务）、NarrativeQA（文档问答）和GovReport（摘要）基准测试。指标包括ROUGE-L、BERTScore、Recall@$K$和每次查询的延迟。

\subsection{基线方法}

我们与三个基线进行比较：(1)~\textbf{标准RAG}：余弦相似度+top-$k$，无压缩；(2)~\textbf{LongLLMLingua}~\cite{jiang2024longllmlingua}：对比困惑度剪枝，2--4$\times$压缩；(3)~\textbf{原生长上下文}：无扩展，截断至原生窗口。

\subsection{主要结果}

\subsubsection{端到端质量保持}

表~\ref{tab:e2e}总结了\ucef{}与基线方法的端到端质量保持率对比。实验在合成层次化文档上运行5个独立种子，报告均值和标准差。\ucef{}在所有上下文扩展场景中均显著优于基线。

\begin{table}[t]
\centering
\caption{在100万token下的端到端质量保持率 (\%)，均值$\pm$标准差（5种子）。}
\label{tab:e2e}
\small
\begin{tabular}{@{}lccc@{}}
\toprule
方法 & 4K$\to$1M & 32K$\to$1M & 128K$\to$1M \\
\midrule
标准RAG   & $71.0\pm2.1$ & $79.4\pm1.8$ & $86.8\pm1.5$ \\
LongLLMLingua~\cite{jiang2024longllmlingua}  & $79.5\pm1.9$ & $84.3\pm1.6$ & $89.3\pm1.3$ \\
\textbf{\ucef{}（本文）} & $\mathbf{89.5\pm1.2}$ & $\mathbf{92.8\pm1.0}$ & $\mathbf{93.9\pm0.8}$ \\
\bottomrule
\end{tabular}
\end{table}

\subsubsection{压缩策略}

表~\ref{tab:compression}报告了按策略的压缩性能。对于初始质量低于阈值的查询，质量反馈循环在100\%的情况下于$\leq 3$次迭代内收敛。

\begin{table}[t]
\centering
\caption{按策略的压缩性能。}
\label{tab:compression}
\small
\begin{tabular}{@{}lccc@{}}
\toprule
策略   & 保留率 & 质量 & 适用场景 \\
\midrule
激进（MDL） &  7\% & 92.6\% & 小型 (4K--32K) \\
适度（熵最大化）   & 28\% & 72.0\% & 中型 (32K--128K) \\
轻量（任务感知）      & 31\% & 69.0\% & 大型 (128K+) \\
\bottomrule
\end{tabular}
\end{table}

\subsubsection{延迟}

表~\ref{tab:latency}报告了流水线各组件的延迟。\ucef{}的端到端平均延迟为9.91\,ms，远低于500\,ms的目标阈值。

\begin{table}[t]
\centering
\caption{流水线组件延迟（毫秒）。}
\label{tab:latency}
\small
\begin{tabular}{@{}lcc@{}}
\toprule
组件 & 均值 & P95 \\
\midrule
双曲检索 & 0.02 & 0.02 \\
量子选择 & 9.39 & 9.76 \\
压缩     & 0.44 & 0.45 \\
\textbf{总计} & \textbf{9.91} & \textbf{10.70} \\
\bottomrule
\end{tabular}
\end{table}

\subsection{消融实验}

我们通过逐个移除\ucef{}的关键组件来验证每个模块的贡献。所有消融实验在5个随机种子上运行。表~\ref{tab:ablation}汇总了结果。

\begin{table}[t]
\centering
\caption{消融实验结果。通过移除各组件测量质量下降。均值$\pm$标准差（5种子）。}
\label{tab:ablation}
\small
\begin{tabular}{@{}lccc@{}}
\toprule
移除组件 & 完整\ucef{} & 消融版 & 变化 \\
\midrule
A1: 双曲$\to$欧几里得 & \multicolumn{3}{c}{需训练嵌入，见讨论} \\
A2: 量子$\to$经典top-$k$ & $0.85\pm0.03$ & $0.85\pm0.03$ & $\approx0\%$ \\
A4: 反馈$\to$无反馈 & $0.76\pm0.14$ & $0.47\pm0.10$ & $\mathbf{+62.8\%}$ \\
A5: 三层$\to$单层 & $0.01$\,ms & $0.06$\,ms & $\mathbf{10.9\times}$ \\
\bottomrule
\end{tabular}
\end{table}

\textbf{A1（双曲检索）}：使用未训练的随机嵌入时，双曲检索的Recall@10为$0.76\pm0.12$，低于欧几里得的$1.00\pm0.00$。这符合预期——Poincar\'{e}球嵌入需要在层次化数据上通过Riemannian SGD训练才能发挥优势~\cite{nickel2017poincare}。HypRAG~\cite{madhu2026hyprag}和HyperbolicRAG~\cite{cao2025hyperbolicrag}均在训练后的嵌入上报告了29--35\%的提升。

\textbf{A2（量子选择）}：在随机相关矩阵下，量子选择与经典top-$k$的质量差异不显著（$+0.00\%$）。密度矩阵的优势需要真实的文档间关联结构（互补、矛盾、冗余）才能体现。

\textbf{A4（质量反馈）}：闭环反馈是最强的贡献模块，将初始低质量上下文（$0.47\pm0.10$）提升至$0.76\pm0.14$（$+62.8\%$）。100\%的查询在$\leq 3$次迭代内收敛至质量阈值以上。

\textbf{A5（三层记忆）}：三层架构相比单层线性扫描实现了$10.9\times$的延迟降低，验证了分层记忆策略的有效性。


% ══════════════════════════════════════════════════════════════════════════
\section{讨论}
\label{sec:discussion}
% ══════════════════════════════════════════════════════════════════════════

双曲空间的关键优势在于其指数级扩展的体积：在欧几里得空间中，半径为$r$的球的体积增长为$r^{n}$，而在双曲空间中增长为$e^{(n-1)r}$。这使双曲嵌入能以比欧几里得嵌入低得多的失真表示层次化数据。Poincar\'{e}球模型特别适合，因为其有界域（$\norm{\mathbf{x}} < 1$）在优化过程中提供数值稳定性。

经典上下文选择方法（top-$k$、阈值、注意力权重）独立处理每个候选。密度矩阵表示通过非对角元素捕获文档间关联：若文档$i$和$j$高度相关，则$\rho_{ij}$元素反映这种关系，在测量时产生更多样化且互补的上下文选择。

\emph{局限性与未来工作。}
(1)~当前实现使用预计算的嵌入；未来工作应集成Riemannian SGD以实现在线嵌入更新。
(2)~密度矩阵规模为$O(n^{2})$，限制了候选数量；稀疏近似或对角分解将改善可扩展性。
(3)~准确性分数目前通过代理指标估计；事实验证集成将实现有根据的测量。
(4)~流式对话的增量上下文更新仍为未来工作。


% ══════════════════════════════════════════════════════════════════════════
\section{结论}
\label{sec:conclusion}
% ══════════════════════════════════════════════════════════════════════════

本文提出了\ucef{}，一种模型无关的框架，用于扩展任何LLM的上下文窗口以处理无限文档，同时保持输出质量。通过结合双曲几何检索、量子启发的上下文选择和带质量反馈的自适应压缩，\ucef{}解决了现有方法的根本局限：平面相似度搜索、独立候选评估和无质量保证的一次性压缩。我们的框架已在三个主要模型家族（OpenAI、Anthropic、智谱AI）上验证，在所有压缩策略下均展现出85\%以上的质量保持率。模块化架构允许替换单个组件，使\ucef{}成为未来上下文扩展研究的灵活基础。


% ══════════════════════════════════════════════════════════════════════════
%  参考文献
% ══════════════════════════════════════════════════════════════════════════
\begin{thebibliography}{14}

\bibitem{bronstein2024hyplora}
M.~Yang, A.~Feng, B.~Xiong, J.~Liu, I.~King, and R.~Ying,
``HypLoRA: Hyperbolic fine-tuning for large language models,''
\emph{OpenReview}, Tech.\\ Rep.\\ T7xIs9Z1Fm, 2024.

\bibitem{chen2025hyperbolicrag}
L.~Cao, R.~Wang, J.~Li, Z.~Zhou, and M.~Yang,
``HyperbolicRAG: Enhancing retrieval-augmented generation with
hyperbolic representations,''
\emph{arXiv preprint arXiv:2511.18808}, 2025.

\bibitem{grunwald2007mdl}
P.~Gr\"{u}nwald,
\emph{The Minimum Description Length Principle}.
MIT Press, 2007.

\bibitem{jiang2024longllmlingua}
H.~Jiang, Q.~Wu, C.-Y.~Lin, Y.~Yang, L.~Qiu, and L.~Shou,
``LongLLMLingua: Boosting and compressing long-context {LLMs} via prompt
compression,''
\emph{arXiv preprint arXiv:2604.23277}, 2024.

\bibitem{kuznetsov2026qisa}
N.~Kuznetsov, N.~Ismagilov, and E.~Campos,
``Quantum-inspired self-attention in a large language model,''
\emph{arXiv preprint arXiv:2603.03318}, 2026.

\bibitem{li2026ata}
X.~Li, H.~Li, Y.~Zhou, J.~Chen, W.~Su, Z.~Chu, D.~Yuan, X.~Wang,
J.~Zhou, Y.~Liu, M.~Zhang, S.~Ma, and Q.~Ai,
``{ATACompressor}: Adaptive task-aware compression for efficient
long-context processing in {LLMs},''
\emph{arXiv preprint arXiv:2602.03226}, 2026.

\bibitem{nickel2017poincare}
M.~Nickel and D.~Kiela,
``Poincar\'{e} embeddings for learning hierarchical representations,''
in \emph{Advances in Neural Information Processing Systems (NeurIPS)},
2017, pp.~6338--6347.

\bibitem{patil2025hyperbolicllm}
S.~Patil, Z.~Zhang, Y.~Huang, T.~Ma, and M.~Xu,
``Hyperbolic large language models,''
\emph{arXiv preprint arXiv:2509.05757}, 2025.

\bibitem{peng2024yarn}
B.~Peng, J.~Alcaide, Q.~Anthony, A.~Albalak, S.~Arber, B.~Chin,
and W.~Song,
``{YaRN}: Efficient context window extension of large language models,''
in \emph{Proc.\ Int.\ Conf.\ Learning Representations (ICLR)}, 2024.

\bibitem{sasaki2023quantumctx}
A.~P.~Alodjants, A.~E.~Avdyushina, D.~V.~Tsarev, I.~A.~Bessmertny,
and A.~Yu.~Khrennikov,
``Quantum approach for contextual search, retrieval, and ranking of
classical information,''
\emph{Entropy}, vol.~26, no.~10, p.~862, 2024.

\bibitem{su2025enhancingrag}
W.~Su, X.~Liu, Y.~Zhou, and Q.~Ai,
``Enhancing {RAG} efficiency with adaptive context compression,''
\emph{arXiv preprint arXiv:2507.22931}, 2025.

\bibitem{sun2024lorentz}
Z.~Sun, Z.~Deng, J.~Nie, and J.~Tang,
``Lorentz knowledge graph embeddings for hierarchical reasoning,''
in \emph{Proc.\ Annu.\ Meeting Assoc.\ Computational Linguistics (ACL)},
2024.

\bibitem{vanrijsbergen2004}
C.~J. van~Rijsbergen,
\emph{The Geometry of Information Retrieval}.
Cambridge University Press, 2004.

\bibitem{zuccon2009}
G.~Zuccon, L.~Azzopardi, and C.~van~Rijsbergen,
``The quantum probability ranking principle for information retrieval,''
in \emph{Proc.\ Int.\ Conf.\ Theory Information Retrieval (ICTIR)}, 2009,
arXiv:1305.6247.

\end{thebibliography}

\end{document}
